\section{Proof of multiplicative adequacy}
\label{sec:annex-p08-imll-complete}

\begin{annexitem}{annex:p08-imll-complete}{thm:p08-imll-complete}{Soundness completeness and cut elimination}
\begin{proof}
We use only the rules of \hyperref[sheet:imll]{the registered IMLL calculus}. A context is a multiset
of occurrences, not a set of formulas. In a cut, one specified occurrence
in the right antecedent is removed. Other occurrences of the same
formula are retained.

\emph{Soundness.}
Fix a commutative residuated partially ordered monoid and a valuation.
Let \(g\) and \(d\) be the products assigned to the two contexts.
Write \(a,b,c\) for the values of \(A,B,C\).
Identity is sound because \(a\leq a\).
For cut, \(g\leq a\) and \(d\otimes a\leq b\) imply
\(d\otimes g\leq d\otimes a\leq b\) by monotonicity and transitivity.
Tensor left changes no interpreted product:
\(g\otimes a\otimes b=g\otimes(a\otimes b)\).
Tensor right follows from
\(g\leq a\), \(d\leq b\), and
\(g\otimes d\leq a\otimes b\).
Implication right is precisely residuation:
\(g\otimes a\leq b\) implies \(g\leq a\multimap b\).
For implication left, its premises give \(g\leq a\) and
\(d\otimes b\leq c\).
Residuation applied to reflexivity
\(a\multimap b\leq a\multimap b\) gives
\(a\otimes(a\multimap b)\leq b\).
Consequently
\[
 \begin{aligned}
 d\otimes g\otimes(a\multimap b)
 &\leq d\otimes a\otimes(a\multimap b)\\
 &\leq d\otimes b\leq c.
 \end{aligned}
\]
Unit right is \(1\leq1\), and unit left changes \(g\) to
\(g\otimes1=g\). Multiset permutation is sound by commutativity.
Induction on derivations proves soundness, including derivations with cut.

\emph{Admissibility of a cut with cut-free premises.}
Suppose cut-free derivations end in
\(\Gamma\Rightarrow A\) and \(\Delta,A\Rightarrow C\).
Mark the displayed occurrence of \(A\) in the second sequent.
We construct a cut-free derivation of \(\Gamma,\Delta\Rightarrow C\).
Induct lexicographically on the number of connective and unit occurrences
in \(A\), followed by the sum of the two derivation heights.

If the left derivation is identity, its context consists of \(A\),
and the right derivation already has the required conclusion.
If the right derivation is identity, its antecedent consists of its
marked occurrence, \(C=A\), and the left derivation suffices.
Unit right cannot end the right premise, since that premise has a
marked antecedent occurrence.

Suppose the last left rule does not introduce its succedent \(A\).
There are exactly three possibilities: unit left, tensor left, or
implication left. For unit left, perform the cut on the premise
before the unit is inserted, then insert that unit by unit left.
For tensor left, perform the cut before its antecedent tensor is
assembled, then apply tensor left to the same two occurrences.
For implication left its premises have the forms
\(\Pi\Rightarrow B\) and \(\Sigma,D\Rightarrow A\), with conclusion
\(\Pi,\Sigma,B\multimap D\Rightarrow A\).
Cut the second premise with \(\Delta,A\Rightarrow C\);
the induction hypothesis gives
\(\Sigma,D,\Delta\Rightarrow C\).
Reapply implication left with the unchanged
\(\Pi\Rightarrow B\). The result is
\(\Pi,\Sigma,\Delta,B\multimap D\Rightarrow C\).
Each recursive cut has the same cut formula and a smaller height sum.

We may therefore assume that the last left rule introduces \(A\).
If the last right rule does not introduce the marked antecedent
occurrence, commute the cut upward on the right as follows.
For unit left or tensor left on other occurrences, cut into its
single premise and reapply that same rule.
For implication right, cut into its premise and then discharge the
same extra antecedent occurrence as before.
For tensor right, the marked occurrence belongs to exactly one
of the two premise contexts. Cut into that premise, leave the other
premise unchanged, and reapply tensor right with the resulting split.
For implication left on an unmarked formula \(B\multimap D\),
the marked occurrence lies either in its argument premise
\(\Pi\Rightarrow B\) or in its continuation premise
\(\Sigma,D\Rightarrow C\).
In the first case cut into the argument premise and keep the
continuation unchanged; in the second cut into the continuation
and keep the argument unchanged. Reapply implication left.
These cases again decrease the height sum and preserve every
unmarked occurrence. They exhaust the nonprincipal right rules.

It remains that both last rules are principal.
If \(A=\mathbf1\), the left rule is unit right, so \(\Gamma\)
is empty. The right rule is unit left with premise
\(\Delta\Rightarrow C\), which is already the desired derivation.
If \(A=B\otimes D\), the left rule has cut-free premises
\(\Gamma_1\Rightarrow B\) and \(\Gamma_2\Rightarrow D\);
the right rule has premise \(\Delta,B,D\Rightarrow C\).
Use the induction hypothesis first to cut \(B\), producing
\(\Gamma_1,\Delta,D\Rightarrow C\), and then to cut \(D\),
producing \(\Gamma_1,\Gamma_2,\Delta\Rightarrow C\).
Both cut formulas are strictly smaller than \(A\).
If \(A=B\multimap D\), the left rule has premise
\(\Gamma,B\Rightarrow D\). The right rule has premises
\(\Pi\Rightarrow B\) and \(\Sigma,D\Rightarrow C\).
Cut \(B\) between \(\Pi\Rightarrow B\) and
\(\Gamma,B\Rightarrow D\), giving
\(\Pi,\Gamma\Rightarrow D\). Then cut \(D\) with
\(\Sigma,D\Rightarrow C\), giving
\(\Pi,\Gamma,\Sigma\Rightarrow C\).
Again both formulas are strictly smaller.
An atomic cut formula has no principal introduction rule on the left
other than identity, already handled.
This completes the lexicographic induction.

Induction on an arbitrary derivation now eliminates every cut.
First replace the immediate subderivations by their cut-free versions.
Reapply a non-cut last rule; for a cut last rule apply the construction
just proved. Thus conditions 1 and 2 of the theorem are equivalent.

\emph{The algebra of formulas.}
For this construction, derivability initially allows cut.
Declare two formulas equivalent when each derives the other from a
singleton antecedent. Identity gives reflexivity, symmetry is built
into the declaration, and cut gives transitivity. Let \([A]\) denote
the equivalence class of \(A\), and put
\[
 [A]\leq[B]\quad\Longleftrightarrow\quad
 A\Rightarrow B\text{ is derivable}.
\]
This is independent of representatives: replace either representative
by cutting with the corresponding derivation of mutual equivalence.
It is a partial order, since antisymmetry is exactly the equivalence
used to form the quotient.

Define the product of \([A]\) and \([B]\) to be \([A\otimes B]\)
and the unit to be \([\mathbf1]\).
If \(A\Rightarrow A'\) and \(B\Rightarrow B'\), tensor right gives
\(A,B\Rightarrow A'\otimes B'\), and tensor left gives
\(A\otimes B\Rightarrow A'\otimes B'\).
Thus product is monotone and well-defined on equivalence classes.
For commutativity, identities and tensor right give
\(A,B\Rightarrow B\otimes A\); tensor left gives
\(A\otimes B\Rightarrow B\otimes A\), and the reverse is symmetric.
For associativity, identities and two tensor-right steps derive
\(A,B,C\Rightarrow A\otimes(B\otimes C)\).
Two tensor-left steps give
\((A\otimes B)\otimes C\Rightarrow A\otimes(B\otimes C)\).
Reversing the target bracketing gives the converse derivation.
For the unit law, identity followed by unit left and tensor left
gives \(A\otimes\mathbf1\Rightarrow A\).
Identity, unit right and tensor right give
\(A\Rightarrow A\otimes\mathbf1\).
These establish a commutative ordered monoid.

Set the residual of \([A]\) into \([B]\) to be
\([A\multimap B]\).
To check representatives, suppose \(A'\Rightarrow A\) and
\(B\Rightarrow B'\). Implication left on these two derivations gives
\(A',A\multimap B\Rightarrow B'\); implication right gives
\(A\multimap B\Rightarrow A'\multimap B'\).
Applying this in both directions proves well-definedness under
mutual equivalence of both arguments.

Finally,
\[
 A\otimes B\Rightarrow C
 \quad\Longleftrightarrow\quad
 A,B\Rightarrow C
 \quad\Longleftrightarrow\quad
 B\Rightarrow A\multimap C.
\]
For the first equivalence, tensor left gives the forward
assembly from the middle sequent. Conversely, tensor right applied
to identities gives \(A,B\Rightarrow A\otimes B\), which cuts
with the left sequent. For the second equivalence, implication
right gives the implication from the middle to the right.
Conversely, implication left on identities gives
\(A,A\multimap C\Rightarrow C\), and cutting the right sequent
into its implication occurrence gives \(A,B\Rightarrow C\).
This proves precisely the required residual law in the quotient.

\emph{Completeness.}
Interpret each atom \(p\) in the quotient monoid as \([p]\).
Induction on formulas gives the value \([A]\) for every formula \(A\).
For a nonempty enumeration of \(\Gamma\), choose any bracketing of
its tensor formula; write it as \(T\). If \(\Gamma\) is empty,
take \(T=\mathbf1\).
Its class is the product of the values in \(\Gamma\), independently
of enumeration and bracketing by the monoid laws just proved.
Identities and tensor right derive \(\Gamma\Rightarrow T\)
in the nonempty case; unit right does so in the empty case.

If condition 3 holds, it holds in this particular quotient monoid.
Hence \([T]\leq[A]\), which means that \(T\Rightarrow A\) is
derivable. Cut with \(\Gamma\Rightarrow T\) gives
\(\Gamma\Rightarrow A\). Thus condition 3 implies condition 1.
Soundness gave the converse, and cut elimination gives condition 2.
If no derivation exists, the same quotient valuation fails the
required inequality, yielding the asserted countermodel.
\end{proof}
\end{annexitem}

\begin{annexitem}{annex:p08-imll-proof-bound}{thm:p08-imll-proof-bound}{The size accounting for cut-free proofs}
\begin{proof}
Give each formula its syntax-tree size, counting each atom, connective
and unit once, and give a sequent the sum of its formula sizes,
including every antecedent occurrence and its succedent.
For each nonaxiom cut-free rule, the sum of the sizes of its premises
is exactly the size of its conclusion minus one.
Tensor left replaces one tensor occurrence by its two factors.
Tensor right distributes the context between its premises and replaces
the succedent tensor by its factors. Implication right exposes the
domain as one antecedent occurrence and retains the codomain as
succedent. Implication left distributes the context and exposes its
domain in one premise and its codomain in the other. Unit left
removes one unit occurrence. These are all nonaxiom rules; none
duplicates a context.

Summing this identity over the internal nodes of a proof tree cancels
the sizes of all nonroot internal sequents. The root size \(N\) is
therefore the number of internal nodes plus the sum of the leaf
sequent sizes. An identity leaf has positive size, as does the
unit-right leaf. Hence the number of internal nodes plus the number
of leaves is at most \(N\). This is the claimed node bound.
The same accounting also shows that every sequent occurring in the
proof has size at most \(N\).

Backward search first tests the two axiom forms. Otherwise it enumerates
each possible last rule, each principal formula occurrence, and every
multiset split required by a binary rule. There are finitely many
choices. In every recursive premise the size is strictly smaller
than that of the current conclusion, so search terminates. All leaves
are checked axioms, and each reconstructed step is a registered rule,
making success sound. Cut elimination ensures that every derivable
input has a cut-free proof among the explored trees, making failure
decisive.

For the complexity claim, use the actual input string length rather
than treating arbitrarily long atom names as unit-cost data.
Every formula in a cut-free proof is a subformula of an input formula.
The node bound and the sequent-size bound therefore yield a polynomial
encoding of the whole tree, even when formulas are written explicitly
at every node. Rule identifiers, principal-occurrence positions and
context splits also have polynomial total length.
A verifier checks each axiom and each rule by finite syntactic
comparison and multiset accounting. These checks take polynomial
time in the input and certificate lengths.
Guessing such a tree and checking it proves membership in
\(\mathsf{NP}\); exhaustive deterministic search need not be efficient.
\end{proof}
\end{annexitem}

\begin{annexitem}{annex:p08-imll-partition}{thm:p08-imll-partition}{The partition reduction}
\begin{proof}
Suppose first that the indices have been partitioned into triples of
sum \(B\).  Fix one triple \(i,j,k\).  After introducing the right-hand
implication of \(D\), decompose the assumption \(p^{\otimes B}\)
by tensor-left into \(B\) occurrences of \(p\).  Divide these occurrences
into blocks of lengths \(a_i,a_j,a_k\).  In each block, identities and
tensor-right derive \(p^{\otimes a_ell}\), and implication-left with
\(F_ell\) derives one occurrence of \(q\).  Tensor-right joins these
three derivations into \(q^{\otimes3}\), using each item and each
occurrence of \(p\) exactly once.  Reassemble the left tensor and
introduce the right implication.  This derives \(D\) from precisely
\(F_i,F_j,F_k\).  Tensor-right then combines the \(m\) disjoint
triple derivations into the encoded sequent.

For the converse, use cut elimination from
\cref{thm:p08-imll-complete}.  Two applications of its soundness statement
give useful exact resource invariants.  The additive group \(\mathbb Z\),
ordered by equality, is a commutative residuated partially ordered monoid:
the monoid operation is addition, the unit is zero, and the residual of
\(x\) in \(y\) is \(y-x\).  Indeed,
\(x+z=y\) holds exactly when \(z=y-x\).  Interpret \(p\) by \(1\)
and \(q\) by \(0\).  Every item \(F_i\) then has value \(-a_i\).
Consequently no submultiset of items can derive
\(p^{\otimes a_j}\): its antecedent has nonpositive value whereas
the succedent has positive value.

Consider now a cut-free derivation of a tensor of one or more copies
of \(D\), with only item formulas in the antecedent.  An implication-left
last inference would require exactly such an impossible argument premise
for one of those items.  Tensor-left and unit-left do not apply to this
antecedent.  An identity leaf is impossible: no item is syntactically
\(D\) or a tensor of copies of \(D\), since an item's consequent is
\(q\), whereas that of \(D\) is \(q^{\otimes3}\).  If more than one
copy of \(D\) remains in the succedent, the last inference must therefore
be tensor-right.  Repeating this argument along the fixed tensor tree
partitions all antecedent occurrences into \(m\) disjoint groups, each
deriving \(D\).  For \(m=1\) the single group is the original context.

Apply soundness again, now in the equality-ordered additive group
\(\mathbb Z^2\), assigning \(p=(1,0)\) and \(q=(0,1)\).
An item has value \((-a_i,1)\), and \(D\) has value \((-B,3)\).
Thus a group of indices \(J\) deriving \(D\) satisfies
\[
 \left(-\sum_{i\in J}a_i,\lvert J\rvert\right)=(-B,3).
\]
Each group contains exactly three indices and has sum \(B\).  The
groups extracted from tensor-right are consequently the required
3-partition.  This uses equality as the order, not the usual order
on the integers; it is conservation, rather than an inequality estimate.

There are \(\sum_i a_i=mB\) copies of \(p\) in the items and
\(mB\) in the target.  All remaining atoms, connectives, and delimiters
contribute \(O(mB+m)\) symbols.  They can be emitted in that time,
which is polynomial in the unary instance length.  A total reduction
on arbitrary input strings first checks the instance conditions in
polynomial time and sends malformed inputs to \(p\Rightarrow q\),
which is underivable by the same equality-ordered integer model.
\end{proof}
\end{annexitem}

\begin{annexitem}{annex:p08-imll-np-complete}{thm:p08-imll-np-complete}{Complexity}
\begin{proof}
Unary 3-PARTITION is NP-complete by the strong NP-completeness result
\parencite[Theorem~4.4]{GareyJohnson1979}.
\Cref{thm:p08-imll-partition} is a polynomial-time many-one reduction
from that problem to derivability in the registered calculus, proving
NP-hardness.  The polynomial cut-free certificates and their verifier
in \cref{thm:p08-imll-proof-bound} prove membership in NP.  Together
these give NP-completeness.  The reduction uses neither the unit nor
atoms other than \(p,q\).
\end{proof}
\end{annexitem}
